478
8 Differential Calculus in Several Variables
\begin{table}[h]
    \centering
    \caption{8.1}
    \begin{tabular}{|c|c|c|c|}
    \hline
    Temperature, $^\circ\text{C}$ & Frequency, $\%$ & Temperature, $^\circ\text{C}$ & Frequency, $\%$ \\
    \hline
    0 & 39 & 20 & 136 \\
    5 & 54 & 25 & 182 \\
    10 & 74 & 30 & 254 \\
    15 & 100 & & \\
    \hline
    \end{tabular}
\end{table}
    for determining the coefficients $a$ and $b$. Here, following Gauss, we write $[x_k, x_k] :=
    \sum_{i=1}^n x_i x_i, [x_k, 1] := \sum_{i=1}^n x_i \cdot 1 + \ldots + x_n \cdot 1, [x_k, y_k] := \sum_{i=1}^n x_i y_i + \ldots + x_n y_n$, and
    so forth.
    \begin{enumerate}
        \setcounter{enumi}{7}
        \renewcommand{\theenumi}{\alph{enumi}}
        \item b) Write the system of equations for the numbers $a_1, \ldots, a_m, b$ to which the
        least-squares principle leads when Eq. (8.79) is replaced by the relation
        \[
        y = \sum_{i=1}^m a_i x^i + b,
        \]
        (or, more briefly, $y = \sum_{i=1}^m a_i \mathbf{x}^i + b$) between the quantities $x^1, \ldots, x^m$ and $y$.
        \item c) How can the method of least squares be used to find empirical formulas of
        the form
        \[
        y = c x_1^{\alpha_1} \ldots x_n^{\alpha_n}
        \]
        connecting physical quantities $x_1, \ldots, x_m$ with the quantity $y$?
        \item d) (M. Germain) The frequency $R$ of heart contractions was measured at differ-
        ent temperatures $T$ in several dozen specimens of \textit{Nereis diversicolor}. The frequen-
        cies were expressed in percents relative to the contraction frequency at $15\,^\circ\text{C}$. The
        results are given in Table 8.1.
        The dependence of $R$ on $T$ appears to be exponential. Assuming $R = Ae^{bT}$, find
        the values of the constants $A$ and $b$ that best fit the experimental results.
        \item a) Show that in Huygens' problem, studied in Example 5, the function (8.71)
        tends to zero if at least one of the variables $m_1, \ldots, m_n$ tends to infinity.
        \item b) Show that the function (8.71) has a maximum point in $\mathbb{R}^n$ and hence the
        unique critical point of that function in $\mathbb{R}^n$ must be its maximum.
        \item c) Show that the quantity $v$ defined by formula (8.72) is monotonically increas-
        ing as $n$ increases and find its limit as $n \to \infty$.
        \item During so-called exterior disk grinding the grinding tool -- a rapidly rotating
        grinding disk (with an abrasive rim) that acts as a file -- is brought into contact with
        the surface of a circular machine part that is rotating slowly compared with the disk
        (see Fig. 8.3).
        The disk $K$ is gradually pressed against the machine part $D$, causing a layer $H$ of
        metal to be removed, reducing the part to the required size and producing a smooth
        working surface for the device. In the machine where it will be placed this surface
        will usually be a working surface. In order to extend its working life, the metal of the
        machine part is subjected to a preliminary annealing to harden the steel. However,
    \end{enumerate}